\section{Numerical realization is not cutpoint simulation}
\label{sec:semantics}
A numerical simulator must preserve the values in \eqref{eq:law},
whereas a strict-cutpoint simulator need only preserve whether the
answer probability is larger than $1/2$. The following construction
makes the difference explicit in the present command model. Its
stochastic embedding mechanism is classical; it is the quantitative
width theorem that prevents numerical exactification at bounded width.

\begin{proposition}[A constant-width sign simulator]\label{prop:sign-simulator}
For any orthogonal action on $\R^D$, $D\ge2$, there is a machine
with $D+1$ command-layer labels, common command rows, and one common
query decoder for all lengths such that
\begin{equation}\label{eq:attenuation}
 \widetilde\Pp(B=1\mid x,w,j)-\tfrac12
  =\tfrac12\lambda^{N+1}\langle e_j,U_wx\rangle,
       \qquad\lambda=1/(2D).
\end{equation}
It preserves every strict comparison of \eqref{eq:law} with $1/2$,
including equality. If $\lambda^{N+1}\le\rho$, its worst numerical
binary-TV error is at least
$(\rho-\lambda^{N+1})/(2\sqrt D)$.
\end{proposition}
\begin{proof}
Let $v_0,\ldots,v_D$ be a regular simplex on the unit sphere of
$\R^D$, centered at zero. Then
\[
 \sum_sv_s=0,\qquad \sum_sv_sv_s^T=\frac{D+1}{D}I,
 \qquad
 b_s(z)=\frac{1+D\langle v_s,z\rangle}{D+1}
\]
are barycentric coordinates for every $z$ in its convex hull. For
$\norm z\le1/(2D)$ they are all positive and sum to one.
Initialize with weights $b_s(\lambda x)$. For command $a$ use
$T_a(s,r)=b_r(\lambda U_av_s)$; these are stochastic rows, independent
of $N$ and of the epoch. Use the terminal mean $d_j(s)=\langle e_j,v_s\rangle$.
The displayed barycentric identity shows inductively that the represented
mean is $\lambda^{N+1}U_wx$. This proves \eqref{eq:attenuation} and
its sign assertion. For every unit $U_wx$ some coordinate has modulus
at least $1/\sqrt D$; binary total variation is half the difference
of its means. This gives the numerical error bound. Only the label
is stored. A held binary output uses a separate two-label cut.
\end{proof}

\begin{corollary}[Separation of numerical and sign-only width]
\label{cor:semantic-separation}
For every nontrivial expanding irreducible experiment in
Corollary~\ref{cor:irreducible-width}, numerical width is of order
$N^{s(U)/2}$ at the stated fixed positive accuracy, but strict-cutpoint
width in the same reset--command--query interface is at most $D+1$
for all lengths simultaneously. For complex conjugation the numerical
order is $N^{q-1}$, while the sign-only command width is at most $q^2$.
\end{corollary}
\begin{proof}
Combine Proposition~\ref{prop:sign-simulator} with the corresponding
numerical width theorem. Here $D=q^2-1$ in the last assertion.
The number counts command-layer labels, not the total state count of
a language recognizer additionally parsing reset and query syntax.
\end{proof}

\subsection{The current quantum-automata comparison}
Chen--Wu's revised generalized-automaton conversion
\cite[Proposition 3.1]{CWapr} maps a $k$-dimensional real
zero-cutpoint generalized finite automaton to a probabilistic automaton
with $k+1$ states. For its parameters $c,\delta>0$ the conversion has
\begin{equation}\label{eq:CW-normalization}
 f_{\rm PFA}(w)-1/(k+1)=\delta c^{|w|} f_{\rm GFA}(w).
\end{equation}
Their Theorem 3.2 consequently gives $q^2+1$ states for strict-cutpoint
simulation of a $q$-state general one-way QFA. The second version of
that paper improves the larger bound in its original version.
Mixed-state linearization is exact before the positive embedding;
the factor in \eqref{eq:CW-normalization} is the distinction relevant
here. The later paper \cite[Proposition 4 and Theorems 11--12]{CWmay}
uses the related conversion to compare hybrid quantum and classical
language state costs. These are results about cutpoint languages,
not fixed nonvanishing numerical accuracy.

For clarity, the algebra behind \eqref{eq:CW-normalization} is that a
uniform-row matrix $Q$ annihilates a zero-sum coordinate subspace,
while a scaled matrix $cK_a$ carries the desired signed dynamics
there. Matrices $Q+cK_a$ are made stochastic by choosing $c$ small.
After $N$ commands the signed deviation is multiplied by $c^N$.
A length-dependent decoder attempting to divide by this factor need
not lie in $[-1,1]$ and is not a stochastic readout. More generally,
a stochastic post-processing cannot increase total variation.
Allowing redesign at every horizon does not overcome the numerical
lower bound, which already permits that redesign.

A prefix--suffix sign or rank argument alone cannot substitute for
the occupation proof. The matrix of centered target responses has
entries $\langle \rho U_px,U_s^*e_j\rangle$ and ordinary rank at most
$D$ at every split, independently of the horizon. Adding the constant
probability column gives rank at most $D+1$. The prepare--test
sign-rank and shattering statements in \cite{CWmay} control a
threshold relation in this fixed operator space; they do not force
the numerical hidden-width growth proved here. Conversely the present
argument gives no growing lower bound for the sign-only task, as
Proposition~\ref{prop:sign-simulator} demonstrates.

Chen's subsequent work \cite{ChenSep}, including its invariant-algebra
simulation results and Section 7's stabilizer-orbit capacity analysis,
also uses strict-cutpoint behavioral equivalence. Its dimensions and
symmetry constraints belong in this comparison. Orbit dimension as
$\dim G-\dim G_v$, and quantum operator-space linearization, are not
claimed new here. Our invariant first takes the \emph{minimum over all
nonzero directions in each irreducible representation}, then the
\emph{maximum across active constituents}; it governs a growing
fixed-accuracy numerical width as the word length grows. These are
different optimization and asymptotic variables. No assertion of
exhaustive priority clearance is made from this comparison.
